\section{Dataset}
\label{sec:dataset}

GitSkills was built by Destefanis, Graziotin, Vaccargiu, and Ortu from public GitHub repositories during July 2026~\cite{gitskills2027}.
It holds 3{,}797{,}117 skill-file rows from 282{,}200 repositories.
Each row is one occurrence of a file, so a skill copied into several repositories appears once per copy.
We use release version~3, published on 11~September 2026, and we downloaded it once and kept it as a local copy, so later updates to the public collection stay out of the analysis.
The dataset is shared under the Creative Commons Attribution 4.0 International License, which allows sharing and adaptation as long as proper attribution is given.

The collection has four tables.

\begin{itemize}
  \item \emph{Skill files}, with one row per file occurrence.
  \item \emph{Repositories}, with metadata such as star counts and the license string of each repository that hosts a skill.
  \item \emph{Companion files}, which are the files bundled with each representative skill.
  \item \emph{Collection notes}, which record how the files were gathered.
\end{itemize}

Commit history is stored on a fetched skill row rather than in a table of its own.

A skill-file row is identified by its repository name (\textit{repo\_full\_name}) and its path within that repository (\textit{path}).
The fields our study relies on are the file name (\textit{filename}), the location class (\textit{location\_class}), which separates the usual \textit{.claude/skills} folder from other skills directories and from all other places, and a content hash (\textit{file\_sha}) that identical copies share.
A flag (\textit{dedup\_primary}) marks one representative row for each distinct text, and only that row stores the text itself (\textit{content}).
Another flag (\textit{frontmatter\_valid}) records whether the header metadata at the top of the file is valid.
Skill rows join the repository table on the repository name, which supplies the hosting repository's star count and license string.

Companion-file rows belong to a representative skill and link to it through its path.
Each one names a bundled file or folder (\textit{entry\_name}, \textit{entry\_type}), gives its size and hash (\textit{entry\_size}, \textit{entry\_sha}), and, when the text was fetched, stores it in \textit{content} up to 100~KB.
In this paper we call that stored text \emph{companion text}.
The unit of analysis for both questions starts from one distinct skill text, as Section~\ref{sec:method} explains.
